% ---- pg-301 （书页 289） ----

\[
\sqrt{(r+1)F_{r+1,n-r-1}(\alpha)}\sqrt{\widehat{\text{Var}}(\hat{\beta}_i)},
\]
其中 $\widehat{\text{Var}}(\hat{\beta}_i)$ 为 $s^2(\mathbf{Z}'\mathbf{Z})^{-1}$ 对应于 $\hat{\beta}_i$ 的对角元. \hfill$\blacksquare$

上述的置信椭球以极大似然估计 $\hat{\boldsymbol{\beta}}$ 为中心, 其方向和大小均由 $\mathbf{Z}'\mathbf{Z}$ 的特征向量和特征值所决定. 若某个特征根接近于零, 则置信椭球在相应特征向量的方向上会变得很长.

实际工作者往往忽略结论 7.5 中区间估计的 ``联合'' 置信性质. 事实上, 他们用 ``每次一个'' 的 $t$ 值 $t_{n-r-1}(\alpha/2)$ 来代替 $\sqrt{(r+1)F_{r+1,n-r-1}(\alpha)}$, 用区间
\phantomsection\label{eq:7-11}
\begin{equation}
\hat{\beta}_i\pm t_{n-r-1}\left(\frac{\alpha}{2}\right)\sqrt{\widehat{\text{Var}}(\hat{\beta}_i)}
\tag{7-11}
\end{equation}
来搜索重要的预测变量.

\medskip
\noindent\textbf{例 7.4}\quad (用房地产数据拟合回归模型)

\hyperref[tab:7.1]{表 7.1} 中的数据来自威斯康星州密尔沃基一个住宅区的 20 个家庭. 记 $z_1=$ 总居住面积(以百平方英尺为单位), $z_2=$ 评估价值(千美元), 及 $Y=$ 售价(千美元). 用最小二乘法拟合模型
\[
Y_i=\beta_0+\beta_1 z_{i1}+\beta_2 z_{i2}+\epsilon_i
\]
用计算机得到
\[
(\mathbf{Z}'\mathbf{Z})^{-1}=\begin{bmatrix}
5.1523&&\\
0.2544&0.0512&\\
-0.1463&-0.0172&0.0067
\end{bmatrix}
\]

\begin{mvtable}
\phantomsection\label{tab:7.1}\mvtabcap{表 7.1\quad 房地产数据}
\begin{tabular}{ccc}
\hline
$z_1$&$z_2$&$Y$\\
总居住面积(千平方英尺)&估价(千美元)&售价(千美元)\\
\hline
15.31&57.3&74.8\\
15.20&63.8&74.0\\
16.25&65.4&72.9\\
14.33&57.0&70.0\\
14.57&63.8&74.9\\
17.33&63.2&76.0\\
14.48&60.2&72.0\\
14.91&57.7&73.5\\
15.25&56.4&74.5\\
13.89&55.6&73.5\\
15.18&62.6&71.5\\
14.44&63.4&71.0\\
14.87&60.2&78.9\\
18.63&67.2&86.5\\
15.20&57.1&68.0\\
25.76&89.6&102.0\\
19.05&68.6&84.0\\
15.37&60.1&69.0\\
18.06&66.3&88.0\\
16.35&65.8&76.0\\
\hline
\end{tabular}
\end{mvtable}
